The CHMM family, fitted with common forward--backward recursions, family-specific M-steps, and quantile-based initialization, captures the three symmetric stylized-fact diagnostics on daily US-equity returns at $K^\star = 3$, while leaving an excess-kurtosis gap whose size is set by the choice of heavy-tailed emission family. A $K = 2$ replication on SPY localises the binding channel: the IS KS pass rate falls to $79.0\%$ (against $91.5\%$ for CHMM-N at $K^\star = 3$) while the absolute growth-rate ACF-MAE stays essentially constant ($0.0501$ versus $0.0462$) and the fitted $K = 2$ spectrum carries the entire absolute growth-rate ACF in a single persistent mode ($\lambda_2 = 0.942$). On these data, then, the marginal distribution, not the ACF, limits the low-$K$ Gaussian fit. The mechanism is empirical rather than structural: a three-component Gaussian scale mixture can in principle reach arbitrarily high kurtosis through a small-weight, large-variance component, but the joint ML fit, disciplined by the bulk of the data, does not select such a configuration; what a finite Gaussian mixture structurally cannot supply is a genuinely power-law tail, since it remains asymptotically light-tailed at every parameter setting. This reading does not contradict Ryd\'{e}n et al.~\cite{ryden1998stylized}, whose subseries were winsorised at four sample standard deviations, bringing the kurtosis within reach of a Gaussian mixture; on our raw data the marginal channel binds instead. On the ACF we agree with their in-principle limit: a finite-state HMM produces only finitely many geometric decay modes, we match the absolute growth-rate ACF only within our lag-252 MAE tolerance, and the structural fix for genuinely slow decay is the semi-Markov chain of Bulla and Bulla~\cite{bulla2006stylized}.

The scope limits are explicit. A CHMM fit once in sample does not generalise across regime introductions: the GLD non-equity stress and the W2 (COVID) and W4 (2022 rate-hike) walk-forward folds broke down under static fitting, and no tested refit cadence closed those folds, while quarterly refit did recover the non-stress cross-ticker panel (Table~\ref{tab:validation}). We evaluated only the three symmetric stylized facts; the dynamic leverage effect is not modelled, although the regime means recover part of the static gain/loss asymmetry. The i.i.d.\ bootstrap and the moment-updated-duration HSMM-N beat the CHMM on raw single-window OoS KS, a marginal-axis result only, and no formal privacy guarantee is claimed for the synthetic output.

What the small model buys is inspectability. CHMM-N at $K^\star = 3$ has $K(K-1) + 2K = 12$ free parameters and the shared-$\nu$ CHMM-t has $15$, against the orders-of-magnitude larger QuantGAN generator--critic architecture~\cite{wiese2020quantgan}. The fitted states read directly as return/volatility regimes: on SPY the three Gaussian states land at a calm regime ($\mu = 0.29$, $\sigma = 0.85$ in annualised growth-rate units), an intermediate regime ($0.14$, $1.91$), and a stress regime ($-0.43$, $3.98$), with the shared-$\nu$ Student-$t$ adding one global tail-shape parameter ($\nu = 5.81$) rather than a per-state shape; no external economic-state validation is performed. BIC/CAIC and rolling-origin validation select the state count, while the closed-form eigenvalue argument explains why additional decay modes are unnecessary on this instance. Future work includes asymmetric emission families targeting the leverage effect, explicit-duration sojourn laws that escape finite geometric memory, online change-point detection against the stress folds, and multi-asset composition through a serial-dependence-preserving copula layer. Code, fitted models, and the runner scripts that regenerate every reported table will be released publicly upon publication; daily prices are from commercial vendors and are not redistributed.
